% ==================================================================
% id:       comb-021
% topic:    קומבינטוריקה
% subtopic: עקרון הכפל וסידורים
% solution: full
% ==================================================================
בקניון קטן יש שש חנויות צמודות, הממוספרות מ-1 עד 6. בהן אמורות לקום שתי חנויות הלבשה (אחת לנשים ואחת לגברים), חנות נעליים, בית מרקחת, חנות ספרים וחנות טבע. בעל הקניון מגריל את מיקום החנויות. כמה סידורים אפשריים יש, אם:
\begin{enumerate}
\item אין הגבלות?
\begin{solbox}
סידור של 6 חנויות שונות ב-6 מקומות: $6!=\ans{720}$.
\end{solbox}
\item חנות מספר 1 תהיה בית מרקחת?
\begin{solbox}
את 5 החנויות האחרות מסדרים ב-5 המקומות הנותרים: $5!=\ans{120}$.
\end{solbox}
\item בין שתי חנויות ההלבשה תהיה רק חנות הנעליים?
\hint{חשבו על שתי חנויות ההלבשה וחנות הנעליים שביניהן כעל "בלוק" של 3 מקומות רצופים.}
\begin{solbox}
שתי חנויות ההלבשה וחנות הנעליים שביניהן יוצרות רצף של 3 חנויות. את הרצף אפשר למקם ב-4 מקומות: $(1,2,3)$, $(2,3,4)$, $(3,4,5)$ או $(4,5,6)$. בתוך הרצף אפשר להחליף בין שתי חנויות ההלבשה: $2!$ דרכים. את 3 החנויות האחרות מסדרים ב-3 המקומות הנותרים: $3!$. סך הכול $4\cdot2\cdot6=\ans{48}$.
\end{solbox}
\item בין שתי חנויות ההלבשה יהיו בדיוק שתי חנויות?
\begin{solbox}
המקומות של שתי חנויות ההלבשה הם $(1,4)$, $(2,5)$ או $(3,6)$: 3 אפשרויות, ועוד $2!$ דרכים לקבוע איזו מהן לנשים. את 4 החנויות האחרות מסדרים ב-4 המקומות הנותרים: $4!$. סך הכול $3\cdot2\cdot24=\ans{144}$.
\end{solbox}
\item בית המרקחת וחנות הטבע יהיו צמודות?
\begin{solbox}
לזוג הצמוד יש 5 מקומות אפשריים: $(1,2),\dots,(5,6)$. בתוך הזוג יש $2!$ סדרים, ואת 4 החנויות האחרות מסדרים ב-$4!$ דרכים: $5\cdot2\cdot24=\ans{240}$.
\end{solbox}
\end{enumerate}
